\documentclass[12pt]{article}
\usepackage[utf8]{inputenc}
\usepackage[T1]{fontenc}
\usepackage[margin=1in]{geometry}
\usepackage{setspace}
\usepackage{xcolor}
\usepackage[colorlinks=true, urlcolor=blue!60!black, linkcolor=blue!60!black]{hyperref}
\setlength{\parindent}{2em}
\setlength{\parskip}{1.25em}
\renewcommand{\baselinestretch}{1.0}
\begin{document}
\begin{center}
{\LARGE \textit{Massachusetts v.\ EPA}, 549 U.S.\ 497 (2007)\par}
\vspace{0.8em}
{\large Case Brief --- Four-Page Edition\par}
\end{center}
\vspace{1em}

\section{Memory Jogger}

Massachusetts had standing---bolstered by ``special solicitude'' for a State asserting a statutory procedural right to protect its territory---to challenge EPA's refusal to regulate greenhouse gases from new cars; greenhouse gases are ``air pollutants,'' and EPA's policy-based refusal was arbitrary and capricious (\href{https://supreme.justia.com/cases/federal/us/549/497/}{Justia}).

\section{Detailed Case Facts}

In 1999 nineteen private organizations petitioned EPA to regulate carbon dioxide, methane, nitrous oxide, and hydrofluorocarbons from new motor vehicles under Section 202(a)(1) of the Clean Air Act, which directs standards for ``any air pollutant'' that in the Administrator's ``judgment'' contributes to air pollution that may endanger public health or welfare; ``welfare'' includes effects on ``weather \ldots\ and climate.'' Two earlier EPA General Counsels had concluded that the Act covered $\mathrm{CO_2}$, and a 2001 National Research Council report found that human activity was warming the planet. In 2003 EPA denied the petition, reasoning (1) that the Act does not authorize regulation of global climate change, invoking \textit{FDA v.\ Brown \& Williamson}, and (2) that even if it did, scientific uncertainty, the President's voluntary programs, and foreign-policy concerns about China and India counseled against regulation. Twelve States, local governments, and environmental groups sought review. Uncontested affidavits showed that rising seas were already eroding Massachusetts's state-owned coast and that U.S.\ transportation produced more than 6\% of global $\mathrm{CO_2}$ emissions.

\section{Procedural History}

A fractured D.C.\ Circuit panel denied review. Judge Randolph assumed authority and upheld EPA's ``judgment''; Judge Sentelle found no standing for a harm ``harmful to humanity at large''; Judge Tatel dissented. After rehearing en banc was denied, the Supreme Court granted certiorari and reversed on April 2, 2007.

\section{Judicial Votes}

\textbf{5--4}, reversed and remanded.
\begin{itemize}
\item \textbf{Majority:} Justice Stevens, joined by Justices Kennedy, Souter, Ginsburg, and Breyer.
\item \textbf{Dissent (standing):} Chief Justice Roberts, joined by Justices Scalia, Thomas, and Alito.
\item \textbf{Dissent (merits):} Justice Scalia, joined by the Chief Justice and Justices Thomas and Alito.
\end{itemize}

\section{Holding}

Massachusetts has standing: a State asserting a congressionally conferred procedural right to protect its quasi-sovereign and proprietary interests gets ``special solicitude,'' widely shared climate risk does not defeat a concrete injury, a meaningful contribution to emissions is traceable, and incremental relief satisfies redressability. Greenhouse gases unambiguously fall within the Act's definition of ``air pollutant.'' Denials of rulemaking petitions are reviewable, and EPA may decline to regulate only on grounds rooted in the statute, not extrastatutory policy. The Court did not decide whether EPA must find endangerment on remand.

\section{Analysis of Opinions}

\textbf{Majority (Stevens).} Under \textit{Lujan}'s footnote 7, a litigant with a procedural right need not meet all normal standards of redressability and immediacy. States are ``not normal litigants'' (\textit{Georgia v.\ Tennessee Copper}): having surrendered self-help to the federal government, Massachusetts deserved special solicitude. Its coastal losses were concrete; U.S.\ vehicle emissions were no trivial cause; and agencies may ``whittle away'' at big problems. On the merits, \textit{Heckler v.\ Chaney} did not bar review; the repeated ``any'' in the definition covers greenhouse gases; \textit{Brown \& Williamson} was distinguishable; and EPA's policy reasons were ``divorced from the statutory text.''

\textbf{Roberts (dissenting).} Nothing in the Act or in \textit{Tennessee Copper} relaxes Article~III for States; \textit{Massachusetts v.\ Mellon} bars \textit{parens patriae} suits against the federal government. Global warming is ``harmful to humanity at large''; the projected losses were neither imminent nor reliably measured; and new U.S.\ cars are a fraction of global emissions overwhelmed by growth abroad. The decision was ``SCRAP for a new generation.''

\textbf{Scalia (dissenting).} The Act does not force EPA to form a judgment on demand, and its reasons for deferring deserved \textit{Chevron} deference. ``Air pollutant'' reasonably excludes globally mixed gases; otherwise everything airborne qualifies.

\section{Examples: Future Applications}

\begin{enumerate}
\item \textbf{Same-Side: State Fisheries.} A State owning tidal lands challenges a denied petition to limit discharges warming its shellfish beds; its territorial injury and procedural right track Massachusetts.
\item \textbf{Same-Side: Downwind Ozone.} A downwind State petitions EPA on interstate ozone with data showing harm to its parks and compliance costs.
\item \textbf{Opposite-Side: Private Group Alone.} A nonprofit without a State co-plaintiff, alleging aesthetic harm from future sea-level rise, likely lacks standing under \textit{Lujan}.
\item \textbf{Opposite-Side: Enforcement Priorities.} States suing over immigration priorities based on indirect fiscal costs lack the statutory right and territorial loss present here.
\item \textbf{Fence-Sitter: Lost Revenue.} A State-run utility's direct financial loss may suffice (\textit{Biden v.\ Nebraska}), but general tax-base erosion from consumer choices usually does not.
\end{enumerate}

\section{Critique}

\textbf{Originalist critique.} Calvin Massey and Jonathan Adler argue that ``special solicitude'' has no basis in Article~III's text or history and invited partisan state attorney-general litigation from both sides. Textualists say the Court stretched a localized air statute to global climate, a tension reflected in \textit{UARG}, \textit{West Virginia v.\ EPA}, \textit{Loper Bright}, and EPA's 2026 move to rescind the Endangerment Finding.

\textbf{Progressive and doctrinal views.} Supporters praise the Court for stopping executive nullification of statutory mandates by pretext. Yet even sympathetic scholars note that the Court never defined what ``special solicitude'' does, leaving lower courts without a rule.

\section{Key Quotations}

\begin{quote}
``Given that procedural right and Massachusetts' stake in protecting its quasi-sovereign interests, the Commonwealth is entitled to special solicitude in our standing analysis.''
\end{quote}
\noindent --- Justice Stevens, for the Court, 549 U.S.\ at 520.

\begin{quote}
``On its face, the definition embraces all airborne compounds of whatever stripe, and underscores that intent through the repeated use of the word `any.' Carbon dioxide, methane, nitrous oxide, and hydrofluorocarbons are without a doubt `physical [and] chemical \ldots\ substance[s] which [are] emitted into \ldots\ the ambient air.' The statute is unambiguous.''
\end{quote}
\noindent --- Justice Stevens, for the Court, 549 U.S.\ at 528--529.

\begin{quote}
``Relaxing Article III standing requirements because asserted injuries are pressed by a State, however, has no basis in our jurisprudence, and support for any such `special solicitude' is conspicuously absent from the Court's opinion.''
\end{quote}
\noindent --- Chief Justice Roberts, dissenting, 549 U.S.\ at 536.

\end{document}
